\section{The Arithmetic Tower}\label{tw:tower}
We construct an infinite pro-$2$ extension with prescribed local groups.
Its finite subextensions will supply the fields in the geometric and
analytic arguments. All pro-$2$ subgroups and normal closures are topological. For a pro-$2$
group $H$, write $D_nH$ for its Zassenhaus filtration, characterized in a
finite quotient by $D_nH=\{h:h-1\in I_H^n\}$, where $I_H$ is the
augmentation ideal of $\mathbb F_2[H]$. The associated restricted Lie
algebra has bracket induced by commutators and restricted square induced
by squaring. Our commutator convention is $[x,y]=x^{-1}y^{-1}xy$.
Put
\[
 T=\{3,5,7,11,13\},\qquad \Sigma=T\cup\{2\},\qquad
 E=\mathbb Q(i,\sqrt2,\sqrt3,\sqrt5,\sqrt7,\sqrt{11},\sqrt{13}).
\]
Let $\mathcal G$ be the Galois group of the maximal pro-$2$ extension of
$\mathbb Q$ unramified at finite primes outside $\Sigma$, with no
restriction at infinity. Fix actual complex conjugation $c$.

\subsection{The Complete Global Presentation}
We use the presentation and cohomological framework for pro-$p$
extensions described in \cite{Koch2002,NeukirchSchmidtWingberg2008}.
The odd local relations and the relation at infinity give a complete
presentation of the global group.
\begin{lemma}\label{tw:complete-global-presentation}
The group $\mathcal G$ has a minimal presentation with $7$
generators and $6$ relators.  The relators can be chosen to be $c^2$ and
the five odd tame local relators.  In particular, the image in the global
free group of any dyadic local relation belongs to their closed normal
closure.
\end{lemma}

\begin{proof}
Set $X=\operatorname{Spec}\mathbb Z[1/\prod_{p\in\Sigma}p]$.
Since $2$ is invertible on $X$, the sheaf $\mu_2$ is constant and Kummer
theory gives
\[
 H^1(X,\mathbb F_2)=
 \langle-1,2,q:q\in T\rangle_{\mathbb F_2}.
\]
The displayed squareclasses are independent, so the generator rank is $7$.

To bound the relation rank, take the inverse system of connected finite
Galois $2$-covers $Y\to X$.
Every class in $H^1(Y,\mathbb F_2)$ is a quadratic torsor.  Its normal
closure over $X$ is again a finite Galois $2$-cover: its Galois group embeds
in an extension of $\operatorname{Gal}(Y/X)$ by a product of copies of
$C_2$.  The direct limit of these first cohomology groups is thus zero. The
five-term sequence of the continuous Cartan--Leray spectral sequence
gives an injection
\begin{equation}\label{tw:global-h2-injection}
 H^2(\mathcal G,\mathbb F_2)
 \lhook\joinrel\longrightarrow H^2(X,\mu_2).
\end{equation}
We obtain the continuous sequence from finite Galois covers, using
the spectral sequence for schemes
\cite[Theorem~14.9]{MilneEtale}.

The Kummer sequence and $\operatorname{Pic}(X)=0$ identify the group on the
right with $\operatorname{Br}(X)[2]$.  Since $X$ is a regular integral
arithmetic curve with finite residue fields, the Brauer residue sequence
identifies $\operatorname{Br}(X)$ with the classes in
$\operatorname{Br}(\mathbb Q)$ having zero invariant outside
$\Sigma\cup\{\infty\}$ \cite[Theorem~6.8.3 and
Example~6.9.4]{Poonen}. Global Brauer reciprocity then identifies its
$2$-torsion with the vectors of local $2$-torsion invariants supported
on $\Sigma\cup\{\infty\}$ whose sum is zero.  In particular its dimension
is $6$. The odd-place and infinite-place coordinates are free, and
reciprocity determines the dyadic coordinate
\cite[Chapter~IV, Example~2.14(e)]{MilneCFT}.

To prove surjectivity, we use the quaternion classes obtained by taking
cups of global quadratic characters.
The class $(-1,-1)$ supplies the infinite coordinate.  If $q\equiv3\pmod4$,
then $(-1,q)$ supplies the coordinate at $q$ among the odd places and
has zero infinite coordinate.  If $(2/q)=-1$, the class $(2,q)$ does the
same.  The remaining primes of $T$ are $5,13$, and $2$ is a quadratic
nonresidue at both. These classes span the odd and infinite coordinates. Their only other
possible nonzero invariant is at two, whose value is determined by
reciprocity.

It follows, by naturality of restriction, that
\begin{equation}\label{tw:local-h2-isomorphism}
 H^2(\mathcal G,\mathbb F_2)
 \xrightarrow{\ \sim\ }
 \bigoplus_{q\in T}H^2(\mathcal G_q,\mathbb F_2)
 \oplus H^2(C_2,\mathbb F_2),
\end{equation}
where $\mathcal G_q$ is the maximal pro-$2$ local group.
At odd $q$, this group has the tame presentation
$\langle\tau_q,\phi_q\mid
\phi_q\tau_q\phi_q^{-1}=\tau_q^q\rangle$ and a one-dimensional second
cohomology group.  The real group is $C_2$ and also has one relation.

For a minimal free pro-$2$ presentation
$1\to R\to\mathcal F\to\mathcal G\to1$, transgression is the
natural perfect pairing
\[
 H^2(\mathcal G,\mathbb F_2)
 \simeq\operatorname{Hom}_{\mathrm{cts}}
       (R/R^2[R,\mathcal F],\mathbb F_2).
\]
This follows from the five-term sequence since $\mathcal F$ is free
\cite[(3.3)--(3.4)]{Quadrelli2021}.
Lift the local generators to $\mathcal F$.  Naturality of transgression
identifies the dual of \eqref{tw:local-h2-isomorphism} with the map taking
each local defining relator to its class in $R/R^2[R,\mathcal F]$.
The odd tame relators and $c^2$ consequently form a basis of this quotient.
The lift of $c$ may be chosen as a free generator, because its image is
nonzero modulo the Frattini subgroup, detected by $\mathbb Q(i)$.

A basis of the relation quotient also gives closed normal generation.
Suppose the indicated relators normally generated $R'\ne R$.
Some finite $2$-quotient of $\mathcal F/R'$ would then have nontrivial image
$N$ of $R/R'$.  The finite nonzero module $N/\Phi(N)$ has nonzero
coinvariants under this finite $2$-group: its group algebra has nilpotent
augmentation ideal.  Hence $N/N^2[N,\mathcal F/R']\ne0$, contradicting
that the indicated relators span $R/R^2[R,\mathcal F]$.
Finally, a dyadic local relation maps to the identity in
$\mathcal G$, so its lift belongs to $R$.
\end{proof}

\subsection{The Prescribed Dyadic Field}
We choose the following extension at two.
\begin{lemma}\label{tw:new-local-two}
There is a Galois extension $L_2/\mathbb Q_2$ with decomposition group
\[
 D=\langle x,y,z\mid x^2,y^2,z^4,[y,z]^2,
                  [x,y],[x,z],[[y,z],z]\rangle
\]
of order $32$, inertia $C_2^3$, residue degree $4$, and normalized
different exponent $9/4$. Its augmentation Hilbert polynomial is
\[
 P_D(t)=(1+t)^3(1+t^2)^2,
\]
and its socle has degree seven.
\end{lemma}
\begin{proof}
Let $U_4/\mathbb Q_2$ be the unramified quartic extension and set
\[
 A=\mathbb Q_2(i,\sqrt5,\sqrt{1+2i}),\qquad
 L_2=A(\sqrt2)U_4.
\]
The norm identity $(1+2i)(1-2i)=5$ gives the dihedral extension
$A/\mathbb Q_2$ of order eight. Its maximal abelian subextension is
$\mathbb Q_2(i,\sqrt5)$ and its unramified degree is two. Thus
$A(\sqrt2)$ has group $D_8\times C_2$, degree sixteen, and unramified
degree two: its elementary abelian abelianization cannot have a cyclic
quotient of order four. Adjoining $U_4$ doubles the degree and the
unramified degree. Its inertia has order eight and is elementary abelian.

Write $w=[y,z]$. The inertia and Frobenius generators can be chosen with
the indicated elementary reciprocity classes $5,-1,-2$. In particular
$z$ is a lift of a uniformizer with the specified unit component.
The elements $x,w$ are central involutions, $y$ is an involution, and
$z$ has order four with $z^2$ independent of $w$. These facts give the
displayed presentation. Conversely, that presentation makes $w$ central:
the identity $w^y=w^{-1}$ and the imposed $w^2=1$ give commutation with
$y$, while commutation with $z$ is imposed explicitly. Its normal forms
$x^a y^b z^c w^d$, with $a,b,d\in\mathbb F_2$ and
$c\in\mathbb Z/4$, have multiplication
\[
 (a,b,c,d)(A,B,C,D)=(a+A,b+B,c+C,d+D+cB).
\]
These normal forms give precisely thirty-two elements. The two
restricted layers have dimensions three and two, which gives the
asserted Hilbert polynomial.

For the different, put $k=\mathbb Q_2(i)$, $\pi=1+i$, and
$\alpha=1+2i=1+\pi^2$. Here $v_k(\pi)=1$ and $v_k(2)=2$.
For $b=1+\pi$, the element
\[
 u=\frac{\sqrt\alpha-b}{\pi}
\]
satisfies the Eisenstein polynomial
\[
 u^2+\frac{2b}{\pi}u+\frac2\pi=0.
\]
It is a uniformizer of $k(\sqrt\alpha)$ and gives an integral power
basis. In its derivative, $2u$ has valuation five and $2b/\pi$ has
valuation two. The relative different has exponent two, and the quadratic
conductor--discriminant formula gives the same exponent for the
character over $k$. Induction to
$\mathbb Q_2$ adds $v_2(\operatorname{disc}k)=2$, so the associated
degree-two representation $\rho$ has conductor exponent four.
For the twist by $\chi_2$, dividing $2\alpha$ by the square $\pi^2$
gives $\alpha'=2-i$. The element $u'=\sqrt{\alpha'}-1$ has Eisenstein
polynomial
\[
 (u')^2+2u'+1-\alpha'=0,
\]
since $v_k(1-\alpha')=1$. Its derivative $2(u'+1)$ has valuation four
in the quadratic extension. Thus the relative conductor is four and
$\rho\otimes\chi_2$ has conductor exponent six. These calculations
use the integral-basis formula for the different and the
conductor--discriminant formula for a quadratic character.

The local genus field $E_2=\mathbb Q_2(i,\sqrt2,\sqrt5)$ has
discriminant exponent sixteen. The conductor--discriminant identity for
$A(\sqrt2)$ gives
\[
 v_2(\operatorname{disc}(A(\sqrt2)))=16+2(4+6)=36.
\]
Its normalized different is $36/16=9/4$, which is unchanged by the
unramified base extension to $L_2$.
\end{proof}

For the choice of generators in this lemma, the elementary reciprocity
classes are $5,-1,-2$, in that order. Here reciprocity classes refer to the
maximal elementary abelian quotient. The characters of $-1$ and $5$
on these generators are respectively $y+z$ and $z$. Consequently the
central square form of $A$ is $(y+z)z$, while the unramified quartic
contributes $z^2$. Their independent sum gives the two central directions
$[y,z]$ and $z^2$ in the presentation. Inertia is
$\langle x,y,[y,z]\rangle$. In particular, we may choose the indicated
representatives of the two inertia directions to have order two.

The maximal pro-$2$ group of $\mathbb Q_2$ has three generators and
one relation. Local Kummer theory gives generator rank three.
The pro-$2$ cohomology comparison in the global proof injects its
second cohomology into $\operatorname{Br}(\mathbb Q_2)[2]$, and the
nonzero local Hilbert pairing makes this injection an isomorphism.
The Kummer basis dual to $5,-1,-2$ under that pairing is $2,-5,5$.
Its Hilbert matrix is
\[
 \begin{pmatrix}0&1&1\\1&1&0\\1&0&0\end{pmatrix},
\]
as follows, for odd $u,v$, from
\[
 (2^a u,2^b v)_2
 =(-1)^{(u-1)(v-1)/4+a(v^2-1)/8+b(u^2-1)/8}.
\]
By the relation--cup identity \cite[Proposition~3.2]{Quadrelli2021},
an actual defining relation in these generators has quadratic initial
$y^{[2]}+[x,y]+[x,z]$.

\subsection{The Global Cuts and Their Finite Quotient}
Define $G$ from $\mathcal G$ by imposing the full kernel of the local
map at two to $D$, inertia squares at $3,5,7,11,13$, Frobenius squares
at $3,5$, and Frobenius fourth powers at
$7,11,13,17,19,23,29,31$. At an odd ramified prime choose the Frobenius
lift with zero coordinate on its own genus character. We may make this
choice by multiplying by inertia. After the inertia square is imposed, the
local tame relation makes inertia and Frobenius commute, so the specified
even-power cut is independent of this remaining choice. Every cut
vanishes in $E$, so the resulting extension still contains $E$.

For the finite binary calculations, number the genus basis
$(-1,2,3,5,7,11,13)$ by $0,\ldots,6$ and encode
$v=\sum v_ie_i$ by the integer $\sum 2^iv_i$. For odd $p$, let
\[
 f_p=\sum_{i:a_i\ne p}\frac{1-(a_i/p)}2e_i,
 \qquad (a_0,\ldots,a_6)=(-1,2,3,5,7,11,13),
\]
where $(\cdot/p)$ is the Legendre symbol. The relevant vectors are
\[
\begin{array}{c|rrrrrrrrrr}
 p&3&5&7&11&13&17&19&23&29&31\\\hline
 f_p&43&86&77&83&58&60&71&57&38&101.
\end{array}
\]
The inertia vector at $p\in T$ is $e_{i(p)}$, and the dyadic vectors are
$x=2$, $y=53$, $z=89$. Let $S(v)$ and $B(v,w)$ denote the free
restricted square and bracket in degree two. In the ordered basis
\[
 e_0^{[2]},\ldots,e_6^{[2]},\quad [e_i,e_j]\ (0\le i<j\le6)
\]
with lexicographically ordered pairs,
\[
 S(v)=\sum_i v_i e_i^{[2]}+\sum_{i<j}v_iv_j[e_i,e_j],\qquad
 B(v,w)=S(v+w)+S(v)+S(w).
\]
The six original quadratic initials are
\[
 S(e_0),\qquad
 B(e_{i(p)},f_p)+\frac{p-1}{2}S(e_{i(p)})\quad(p\in T),
\]
with scalars reduced modulo two. These formulas follow by expanding the tame
relators and the infinity
square. We use them to compute the first three layers of the quotient.

\begin{lemma}\label{tw:new-retained-quotient}
The actual quotient $Q=G/D_4G$ has order $2^{38}$. It retains the
entire local group $D$, local groups $C_2^2$ at $3,5$,
$C_2\times C_4$ at $7,11,13$, and $C_4$ at
$17,19,23,29,31$. The conjugacy class of actual complex conjugation
in $Q$ has size $2^{12}$. Its elementary image is outside every one
of these decomposition groups' elementary spans.
\end{lemma}
\begin{proof}
Use the genus basis $(-1,2,3,5,7,11,13)$. The dyadic generators have
elementary vectors
\[
 x=2,\qquad y=53,\qquad z=89.
\]
The six original odd and infinite quadratic relations are independent,
and their sum is $y^{[2]}+[x,y]+[x,z]$.
Lemma~\ref{tw:complete-global-presentation} shows that the corresponding
original relators give a complete presentation.

The local presentation in Lemma~\ref{tw:new-local-two} has four
independent quadratic initials. In its local kernel relation quotient,
the actual defining relation $r$ of the maximal pro-$2$ group of
$\mathbb Q_2$ has coefficient one on $y^2$: its quadratic initial is
$y^{[2]}+[x,y]+[x,z]$. The other six local kernel generators have no
$y^{[2]}$ coordinate. Replacing $y^2$ by $r$ preserves generation of
the relation quotient, hence closed normal generation by the pro-$2$
Nakayama argument. Globally $r$ already belongs to the closed normal
closure of the six original relators. We need not impose it
again.

The resulting complete presentation has sixteen independent quadratic
initials: the six originals, five odd inertia squares,
$x^{[2]},[x,y],[x,z]$, and the two Frobenius squares at $3,5$.
It also has the actual cubic $[[y,z],z]$ and ten degree-four relators,
including both $z^4$ and $[y,z]^2$. Reduction of the local square and
commutator forms against these sixteen rows retains exactly the two
independent local classes $z^{[2]}$ and $[y,z]$. The three local
degree-one generators are independent. Thus $G/D_3G$, of order
$2^{7+12}$, already retains all of $D$. Reduction against the same rows
verifies the other displayed local
groups and the separation of the elementary image of complex conjugation.

The free restricted cubic layer has dimension $112$. Bracketing the
sixteen quadratic initials with the seven generators gives rank $92$.
The actual cubic $[[y,z],z]$ is independent and raises this rank to
$93$. To see that there are no other degree-three initials, work in the free quotient
$\mathcal F/D_4\mathcal F$. Its second Zassenhaus subgroup is elementary
abelian and its third is central. The normal closure of the actual
quadratic words $\rho_i$ and actual cubic $\kappa$ is the vector space
spanned by $\rho_i,[\rho_i,X_k],\kappa$. Independence of the sixteen
quadratic initials implies that a combination of the $\rho_i$ can lie in
degree three only with every coefficient zero. Its intersection with
degree three is precisely the span of the displayed brackets
and $\kappa$. Their initials depend only on the quadratic initials, so
unknown cubic coefficients of the actual quadratic words do not change
the calculation. Consequently
\[
 \dim D_3G/D_4G=19,\qquad |Q|=2^{7+12+19}=2^{38}.
\]
In the preceding 28-coordinate convention, the sixteen rows have
hexadecimal masks
\[
\begin{gathered}
\texttt{1},\ \texttt{142104},\ \texttt{1444000},\ \texttt{4480410},\
\texttt{a010820},\ \texttt{d020000},\\
\texttt{4},\ \texttt{8},\ \texttt{10},\ \texttt{20},\ \texttt{40},\
\texttt{2},\ \texttt{1a080},\ \texttt{2c080},\
\texttt{814aab},\ \texttt{42aa056}.
\end{gathered}
\]
We compute the ranks by binary elimination in a free associative algebra.
Embed the free restricted degree-two space in the free associative algebra
on $X_0,\ldots,X_6$ by
$e_i^{[2]}\mapsto X_i^2$ and
$[e_i,e_j]\mapsto X_iX_j+X_jX_i$.
Represent homogeneous cubics by their 343 coefficients on the words
$X_iX_jX_k$. For each of the sixteen rows $r$ and each $k$, insert
$rX_k+X_kr$ into a binary row-echelon basis. This produces rank 92.
Insert the tensor for $[[y,z],z]$ to obtain rank 93.
Running the same operation on all 28 free quadratic basis elements gives
rank 112. Ordinary binary elimination in degree two gives rank 16 and
verifies independence of the two retained local classes. Reduction of
$B(e_0,e_k)$ modulo the sixteen rows gives rank six and kernel
$\langle e_0\rangle$. Reduction of $rX_0+X_0r$, for the 28 free quadratic
basis elements $r$, modulo the cubic relation space gives rank six.
Each calculation uses only the displayed masks and these tensor
operations. The supplementary code performs the same integer operations.

Commutation with $c$ has rank six in the second layer and rank six on
the second layer with values in the third. Its degree-one kernel is
exactly $\langle c\rangle$. Since $D_2Q$ has order $2^{31}$ and its
commutator image with $c$ has order $2^6$, the intersection of the
centralizer with $D_2Q$ has order $2^{25}$. Its image in $Q/D_2Q$ is
the line through $c$, giving $|C_Q(c)|=2^{26}$ and the asserted
conjugacy-class size.
\end{proof}

\subsection{Filtered Fox Images}
We use left Fox coefficients \cite{Fox1953}:
\[
 w-1=\sum_i\partial_iw\,(x_i-1),\qquad
 \partial_i(uv)=\partial_i u+u\partial_i v.
\]
Derivative basis vectors have degree one. If a finite-dimensional space
$M$ has a decreasing filtration $F^0M=M$, eventually zero, write
$H_M(t)=\sum_n\dim(F^nM/F^{n+1}M)t^n$. Subspaces have the induced
filtration and quotients the quotient filtration. For $0<t<1$,
\[
 H_M(t)=(1-t)\sum_{n\ge1}\dim(M/F^nM)t^{n-1}.
\]
Consequently a filtration-preserving surjection $M\to N$ satisfies
$H_N(t)\le H_M(t)$, even without strictness. In particular, for subspaces
$U,V$ of a filtered ambient space,
\begin{equation}\label{tw:filtered-sum}
 H_{U+V}(t)\le H_U(t)+H_V(t)-H_{U\cap V}(t).
\end{equation}
To obtain this inequality, apply additivity to the quotient filtration
from $U\oplus V$, then compare it with the induced filtration on the sum.
We next compute the cost of a full local relation kernel.

\begin{lemma}\label{tw:filtered-local-freeness}
Let $P$ be a finite $2$-group and $D'\le P$ one of
$C_2,C_4,C_2^2,C_2\times C_4$, or the group $D$ of
Lemma~\ref{tw:new-local-two}. Suppose its degree-one directions, and
its nonzero degree-two directions when present, embed in the corresponding
Zassenhaus layers of $P$. If $D'$ has $r$ minimal generators and
augmentation Hilbert polynomial $P_{D'}$, its full induced local Fox
kernel has Hilbert polynomial
\begin{equation}\label{tw:local-full-cost}
 \left(rt-1+P_{D'}(t)^{-1}\right)H_{\mathbb F_2[P]}(t).
\end{equation}
Its image under a filtration-preserving map has no larger Hilbert value.
\end{lemma}
\begin{proof}
These local groups have no nonzero layer above degree two, so the
hypotheses imply $D'\cap D_nP=D_nD'$ for every $n$.
Extend a homogeneous basis of the graded restricted Lie algebra of $D'$
to one for $P$, ordering the local factors last. The restricted PBW
form of Jennings' theorem \cite{Jennings,Quillen} gives lifts $t_j$ of the
remaining monomials such that
\[
 \mathbb F_2[P]=\bigoplus_jt_j\mathbb F_2[D']
\]
as right modules and filtered vector spaces. The shift of the $j^{th}$
summand is the weight of $t_j$, and the shift polynomial is
$H_{\mathbb F_2[P]}/P_{D'}$.
Put $B=\mathbb F_2[D']$. The local map
$B^r(-1)\to I_{D'}$, $(a_i)\mapsto\sum a_i(d_i-1)$, is strict, because
$I_{D'}^n=\sum_iI_{D'}^{n-1}(d_i-1)$. Its kernel has polynomial
$rtP_{D'}-(P_{D'}-1)$. Filtered right-freeness preserves exactness and
strictness under induction, proving the formula. Finally, the local
generator words have independent elementary images in a minimal global
free presentation and so extend to a free pro-$2$ basis. The Fox
change-of-basis matrix and its inverse preserve augmentation degree.
Thus the identification with a coordinate submodule in the global
derivative module preserves the induced filtration. The assertion about
images follows from the cumulative Hilbert formula.
\end{proof}

The dyadic block shares a relation with the odd and infinite blocks.
We compute the image of this relation as well as the full dyadic cost.
\begin{lemma}\label{tw:new-local-fox}
For this dyadic local group the full induced local Fox-image cost is
$d_D(t)=3t-1+P_D(t)^{-1}$. The common actual dyadic defining row has
normalized image polynomial
\[
 s_D(t)=t^2\left(1-\frac{t^7}{P_D(t)}\right).
\]
\end{lemma}
\begin{proof}
The independent local degree-one and degree-two classes in
Lemma~\ref{tw:new-retained-quotient} give an embedding of the local
restricted graded Lie algebra in the global one. Extend its restricted
PBW basis to a global basis, with the local factors last. This proves
filtered freeness as a right local-algebra module, including the exact
augmentation shifts. The local augmentation sequence then gives the
full Fox-image cost $3t-1+P_D^{-1}$ by
Lemma~\ref{tw:filtered-local-freeness}. This argument does not require
commutativity.

Write $B=\mathbb F_2[D]$. The linear coefficients of the actual
local relation's Fox row span $I_B/I_B^2$, because the local Hilbert
pairing is nonsingular \cite[Chapter~III, Theorem~4.4]{MilneCFT}, and
its matrix is the matrix of linear Fox coefficients
\cite[Proposition~3.2]{Quadrelli2021}. In our coordinates the relation
has initial $y^{[2]}+[x,y]+[x,z]$, whose three derivative coefficients
are $y+z,y+x,x$ and are independent. They generate $I_B$ as a right ideal by
Nakayama. Their common left annihilator is the one-dimensional left socle of $B$.
We can check the required graded assertion in this 32-element algebra.
Starting with its displayed normal-form multiplication,
form $J^0=B$ and obtain $J^{n+1}$ by spanning the right products of a
basis of $J^n$ by $x-1,y-1,z-1$. The dimensions for $0\le n\le8$ are
\[
 32,31,28,23,16,9,4,1,0.
\]
On $J^n/J^{n+1}$ the simultaneous multiplication map to
$(J^{n+1}/J^{n+2})^3$ has kernel zero for $n<7$ and dimension one for
$n=7$. This is a finite row reduction on vectors of length 32 using the
specified multiplication. The top vector is the sum of all group elements.
Changing the three degree-one coefficients by an invertible matrix does
not change these kernels. Thus every nonzero homogeneous element below
degree seven has nonzero product with at least one linear Fox coefficient. Row
multiplication raises coefficient degree by exactly one, and the
derivative basis shift contributes one more. The row image consequently has
polynomial $t^2(P_D-t^7)$. Filtered induction gives the stated normalized
cost.

The actual row belongs to the full dyadic block because the local
relation vanishes in $D$. It also belongs to the sum of the odd blocks
and infinity square. Indeed, by the complete global presentation, the
local relation belongs to the closed normal closure of those original
relators. If a defining relator $\rho$ evaluates to one in the finite
group, the Fox product rule gives
\[
 \partial(u\rho u^{-1})=u\partial\rho.
\]
Products of conjugates give module combinations of the original rows.
Evaluation in the finite derivative module is continuous
and its row span is closed, so the same conclusion holds for the closed
normal closure. The filtered local coordinate change used above preserves
the computed image polynomial. Thus the common submodule is an actual
intersection in the global derivative module.
\end{proof}

\subsection{Infinitude and the Resulting Number Fields}
The local Fox bounds imply that the global group is infinite and give
the following family of number fields.
\begin{theorem}\label{tw:new-field-family}\label{tw:field-family}
The group $G$ is infinite. For every sufficiently large power of two $d$
there is a totally imaginary Galois extension $K/\mathbb Q$ of degree
$2d$ retaining $Q=G/D_4G$. For its actual complex conjugation $c$, the
fixed field $F=K^{\langle c\rangle}$ has at least one real embedding and
\[
 \operatorname{rd}(K)=\operatorname{rd}(F)
 =2^{9/4}\sqrt{15015},\qquad
 \frac{4095}{8192}\le\theta:=\frac{r_2(F)}d<\frac12.
\]
The extension $K/F$ is quadratic, unramified at every finite place, and
split at every place above the following rational primes, whose uniform
local types in both $K$ and $F$ are
\[
\begin{array}{c|c|c}
 p&e_p&f_p\\\hline
 2&8&4\\
 3,5&2&2\\
 7,11,13&2&4\\
 17,19,23,29,31&1&4.
\end{array}
\]
\end{theorem}
\begin{proof}
Suppose first that $G$ is finite and put $A=\mathbb F_2[G]$.
The strict global augmentation map $A^7(-1)\to I_G$ has kernel $L$ with
$H_L=1-(1-7t)H_A$. The continuous Fox relation sequence says that the
rows of the complete defining relations generate $L$. This also follows
by tensoring the free augmentation sequence and using the ideal normally
generated by the defining relators.
The local hypotheses of Lemma~\ref{tw:filtered-local-freeness} hold by
Lemma~\ref{tw:new-retained-quotient}. Write
\[
 c_1=\frac{t^2}{1+t},\qquad c_2=2t-1+(1+t)^{-2},\qquad
 c_{24}=2t-1+\frac1{(1+t)^2(1+t^2)},\qquad
 c_4=\frac{t^4}{(1+t)(1+t^2)}.
\]
The two odd $C_2^2$ blocks cost $2c_2H_A$, the three odd
$C_2\times C_4$ blocks cost $3c_{24}H_A$, infinity costs $c_1H_A$,
and the five unramified order-four caps cost $5c_4H_A$.
The full dyadic block costs $d_DH_A$. Its actual intersection with the
sum of the odd and infinite blocks contains the common row image
$s_DH_A$, which is subtracted once by \eqref{tw:filtered-sum}.
Consequently
\[
 1\le H_A(t)\mathcal P(t),\qquad
 \mathcal P(t)=1-7t+2c_2+3c_{24}+d_D+c_1-s_D+5c_4.
\]
Exact rational substitution gives
\[
 \mathcal P(11/34)=-\frac{160218343}{112275691650}<0.
\]
This contradicts finiteness. The inequality includes every defining
relation, and the subtracted term comes from the common row image in the
global Fox kernel.

An infinite profinite group has finite quotients of unbounded order.
Taking common refinements with $Q$ gives arbitrarily large finite
$2$-quotients $H\twoheadrightarrow Q$. A nontrivial kernel meets $Z(H)$
nontrivially by the class equation for conjugation, and hence contains a
central subgroup of order two. Repeated quotienting by such subgroups
inside the kernel realizes every intermediate power-of-two order while
retaining $Q$. These quotients correspond to the claimed Galois fields.

Because $K$ contains $i$, it is totally imaginary. The chosen complex
conjugation fixes a real subfield in the given embedding, so $F$ has a
real embedding. If $b=r_1(F)$ and $H=\operatorname{Gal}(K/\mathbb Q)$,
embeddings of $F$ correspond to cosets of $\langle c\rangle$ in $H$.
A coset is real exactly when its representative centralizes $c$. Hence
\[
 b=|C_H(c)|/2,\qquad \frac bd=\frac{|C_H(c)|}{|H|}.
\]
The conjugacy class in $H$ maps onto the class in $Q$, of size $2^{12}$.
Thus $0<b/d\le2^{-12}$, giving the stated interval for
$\theta=(1-b/d)/2$.

Conjugation preserves an elementary image. The separation verified in
Lemma~\ref{tw:new-retained-quotient} excludes every conjugate
of $c$ from each ramified or selected decomposition group. Intersection
with $\langle c\rangle$ is trivial, proving relative splitting at these
places. Elsewhere $K/\mathbb Q$ is unramified, so $K/F$ is unramified
there also. Relative splitting transfers the displayed local types to $F$.
Tame inertia of order two contributes $p^{1/2}$ to the root discriminant
at each $p\in T$, and the dyadic lemma contributes $2^{9/4}$. Finally
$\Delta_K=\Delta_F^2$ because the relative finite discriminant is one.
\end{proof}
